\documentclass{article}
\usepackage[UTF8]{ctex}
\usepackage{amsmath,amssymb,amsfonts}
\usepackage{enumitem}
\usepackage{geometry}
\geometry{a4paper, margin=1in}
% === 错题本排版优化（不参与题目标识） ===
\setlength{\parindent}{0pt}
\setlength{\parskip}{0.35em plus 0.10em minus 0.08em}
\setlength{\emergencystretch}{3em}
\setlist{itemsep=0.25em, topsep=0.40em, parsep=0pt}
\allowdisplaybreaks[2]
% === 错题本排版优化结束 ===
\begin{document}
\section*{华中师范大学2026年高等代数真题}
\section*{131．华中师范大学2026年高等代数试卷}
一．填空题。\\

\subsection{题目1}
1．设 \(A\) 是 \(n\) 阶方阵，\(A^{*}\) 是其伴随矩阵，若 \(\operatorname{det}(A)=2\) ，则 \(\operatorname{det}\left(\left(3 A^{*}\right)^{-1}+2 A\right)=\) \(\_\_\_\_\)．

\subsection{题目2}
2．设 \(A\) 为 \(n\) 阶方阵，\(A^{2}=A, \operatorname{rank}(A)=r\) ，则 \(\operatorname{det}\left(A+2 E_{n}\right)=\) \(\_\_\_\_\)．

\subsection{题目3}
3．设 \(A, B\) 为 \(n\) 阶方阵，\(A B-B A=A\) ，则 \(\operatorname{tr}\left(A^{2026}\right)=\) \(\_\_\_\_\)．\\
\(\_\_\_\_\)

\subsection{题目4}
4．设 \(V\) 是数域 \(F\) 上的 5 维向量空间，\(W_{1}\) 和 \(W_{2}\) 都是 \(V\) 的 3 维子空间，则 \(\operatorname{dim}\left(W_{1} \cap W_{2}\right)\) 的可能值为 \(\_\_\_\_\)。

\subsection{题目5}
5．实二次型 \(f\left(x_{1}, x_{2}, x_{3}, x_{4}\right)=x_{1} x_{2}+x_{3}^{2}-x_{3} x_{4}-x_{4}^{2}\) 的正惯性指数为 \(\_\_\_\_\)．

\subsection{题目6}
6．已知三阶正交阵 \(A\) 的迹 \(\operatorname{tr}(A)=-2\) ，则 \(A\) 的特征值为 \(\_\_\_\_\)．
\end{document}
